% Figure: free-body diagram plus kinetic diagram for a cylinder rolling
% down an incline. Left: forces. Right: m a_G and I alpha.
\begin{figure}[htbp]
\centering
\begin{tikzpicture}[
  force/.style={draw=engred, -{Stealth}, very thick},
  kin/.style={draw=engblue, -{Stealth}, very thick},
  surf/.style={draw=engblack, thick},
  lbl/.style={font=\small},
]
% == Left panel: free-body diagram ==
\begin{scope}[shift={(0,0)}]
% incline: angle ~ 25 deg for clarity
\draw[surf] (-0.3,0) -- (4.3,2.15);
\draw[surf] (-0.3,0) -- (-0.3,-0.2);
% angle arc
\draw (1.0,0.0) ++(0:0.9) arc (0:26.6:0.9);
\node[font=\scriptsize] at (1.95,0.22) {$\beta$};
% cylinder centre on incline
\coordinate (C) at (2.2,1.55);
\draw[draw=engblack, thick, fill=enggray!12] (C) circle (0.55);
\node[lbl] at (C) {$G$};
% weight (down)
\draw[force] (C) -- ++(0,-1.4) node[below, lbl] {$m\vec{g}$};
% normal (perpendicular to incline, pointing away)
\draw[force] (C) -- ++(116.6:1.25) node[above left, lbl] {$\vec{N}$};
% friction up the slope (prevents slipping)
\draw[force] (C) -- ++(206.6:1.25) node[below left, lbl] {$\vec{f}$};
\node[font=\scriptsize, engred] at (1.0,2.3) {free-body diagram};
\end{scope}
% == equals sign ==
\node at (5.4,0.9) {\Large $=$};
% == Right panel: kinetic diagram ==
\begin{scope}[shift={(6.0,0)}]
\draw[surf] (-0.3,0) -- (4.3,2.15);
\coordinate (C2) at (2.2,1.55);
\draw[draw=engblack, thick, dashed] (C2) circle (0.55);
\node[lbl] at (C2) {$G$};
% m a_G down the slope
\draw[kin] (C2) -- ++(206.6:1.55) node[below left, lbl] {$m\vec{a}_{G}$};
% I alpha (curved arrow, clockwise rolling)
\draw[kin, -{Stealth}] (C2) ++(40:0.85) arc (40:-150:0.85);
\node[engblue, font=\small] at ($(C2)+(-0.05,1.1)$) {$I_{G}\alpha$};
\node[font=\scriptsize, engblue] at (1.0,2.3) {kinetic diagram};
\end{scope}
\end{tikzpicture}
\caption[Free-body and kinetic diagram for a rolling cylinder]{The
paired free-body diagram (left, forces in red) and kinetic diagram
(right, the inertia terms $m\vec{a}_{G}$ and $I_{G}\alpha$ in blue)
for a solid cylinder rolling without slipping down an incline of
angle $\beta$. Setting the two diagrams equal component by component
is Newton's second law for the rigid body: $\sum\vec{F}=m\vec{a}_{G}$
along and normal to the slope, and $\sum M_{G}=I_{G}\alpha$ about the
mass centre. The friction force $\vec{f}$ provides the moment that
spins the cylinder up; the no-slip constraint $a_{G}=\alpha R$ closes
the system. The kinetic-diagram convention keeps the inertia terms
visibly separate from the forces and prevents the common error of
double-counting $m\vec{a}$ as if it were an applied force.}
\label{fig:vol03ch04:incline-fbd}
\end{figure}
